% Exercises for Chapter 19 (Stochastic volatility). Statements only: solutions are kept privately.
% Unless stated otherwise: the Heston model of the notes under Q, the test case S = K = 100, T = 1,
% r = 1%, y = 2%, and V_0 = theta = 0.04, kappa = 2, xi = 0.3.

\begin{exercise}[One option is enough]\label{ex:sv-risk-premium}
\easy Under $\Prob$ let $\dd S=\mu S\,\dd t+S\sqrt V\,\dd B^S$ and $\dd V=\kappa^{\Prob}(\theta^{\Prob}-V)\,\dd t+\xi\sqrt V\,\dd B$,
with $\dd\langle B^S,B\rangle_t=\rho\,\dd t$, $|\rho|<1$, and a dividend yield $y$. Let $C(t,S,V)$ and $C_1(t,S,V)$ be
smooth prices of two options, with $\partial_VC_1\ne0$, and let $\mathcal L^{\Prob}$ be the generator of $(S,V)$
under $\Prob$ (first and second derivatives, no time derivative).
\begin{enumerate}[label=(\alph*)]
  \item Show that a self-financing portfolio long one $C$, short $\Delta_1$ units of $C_1$ and short $\Delta_S$
    shares (dividends reinvested in cash) is locally riskless for suitable $\Delta_1$ and $\Delta_S$, and deduce that
    \[
      \frac{\partial_tC+\mathcal L^{\Prob}C-(\mu+y-r)S\,\partial_SC-rC}{\xi\sqrt V\,\partial_VC}
    \]
    takes the same value $\lambda(t,S,V)$ for every option. (No-arbitrage: a riskless portfolio earns $r$.)
  \item With $\lambda=\kappa_V\sqrt V$ for a constant $\kappa_V$, show that every option price solves
    \noteref{eq:sv-pde} with the parameters $\kappa=\kappa^{\Prob}+\xi\kappa_V$ and $\theta=\kappa^{\Prob}\theta^{\Prob}/\kappa$.
    Can Feller's condition hold under $\Prob$ and fail under $\Q$?
  \item Why would hedging with the stock and \emph{two} options, as in the course's notes, be redundant?
\end{enumerate}
\end{exercise}

\begin{exercise}[Degenerate cases]\label{ex:sv-degenerate}
\easy You may use \noteref{thm:sv-cf} and \noteref{prop:st-mixing}.
\begin{enumerate}[label=(\alph*)]
  \item Solve \noteref{eq:sv-riccati} for $\xi=0$ and show that the price of any European payoff is its
    Black--Scholes price with total variance $w=\theta T+(V_0-\theta)(1-e^{-\kappa T})/\kappa$. Price the put for
    $V_0=0.09$, $\kappa=1.5$, $\theta=0.04$ (and the test case $S$, $K$, $T$, $r$, $y$). Does $\rho$ matter?
  \item Let $\rho=-1$. Write $S_T$ as an explicit function of $(V_T,Q_T)$. Use it to estimate the at-the-money
    put by Monte Carlo from exact variance paths (you may use \texttt{stratified.cir\_exact} on $64$ steps and the
    trapezoidal rule for $Q_T$), and compare with the Fourier price.
  \item Let $\rho=0$. Show that $\varphi(u)=e^{iu(\ln S+(r-y)T)}\,\E\big[e^{-\frac12(u^2+iu)Q_T}\big]$, and that for
    $\rho=0$ the Riccati equation \noteref{eq:sv-riccati} is the equation of the Laplace transform
    $\E[e^{-\lambda Q_T}]$ with $\lambda=\frac12(u^2+iu)$.
\end{enumerate}
\end{exercise}

\begin{exercise}[The law of the variance]\label{ex:sv-variance-law}
\medium Let $\lambda\ge0$.
\begin{enumerate}[label=(\alph*)]
  \item Look for $\E[e^{-\lambda V_t}]=\exp(A(t)+B(t)V_0)$, write the equations for $A$ and $B$, and solve them:
    \[
      \E[e^{-\lambda V_t}]=(1+2\lambda c_t)^{-2\kappa\theta/\xi^2}\exp\Big(-\frac{\lambda e^{-\kappa t}V_0}{1+2\lambda c_t}\Big),\qquad
      c_t=\frac{\xi^2(1-e^{-\kappa t})}{4\kappa}.
    \]
    Justify that the guess is the Laplace transform (the local martingale you build is bounded).
  \item Deduce that $V_t$ converges in law, as $t\to\infty$, to the gamma law with shape $2\kappa\theta/\xi^2$ and
    scale $\xi^2/(2\kappa)$, and give its mean and variance.
  \item How does the stationary density behave near $0$? Relate this to Feller's condition. Compute the stationary
    probability that $V<0.01$ in the test case and for $\kappa=0.5$, $\theta=0.04$, $\xi=1$.
\end{enumerate}
\end{exercise}

\begin{exercise}[Moment explosions]\label{ex:sv-explosion}
\medium Use the notation of \noteref{thm:sv-cf}(iv).
\begin{enumerate}[label=(\alph*)]
  \item For $a(a-1)>0$, solve the real Riccati equation $\partial_\tau D=\frac12\xi^2D^2-\kappa_aD+\frac12a(a-1)$,
    $D(0)=0$, and find its blow-up time in each case of \noteref{eq:sv-explosion}.
  \item Let $\kappa=0.5$, $\xi=1.2$, $\rho=0.6$. Compute $T^*(2)$ by hand. For $T=2$, what does this imply for the
    variance of the plain Monte Carlo estimator of a call, and for the damping in Carr and Madan's formula?
  \item For $\kappa=0.5$, $\xi=1$, $\rho=-0.9$, and for the Eurostoxx parameters of the notes, compute
    $a_+(T)=\sup\{a\ge1:\E[S_T^a]<\infty\}$ at $T=1$ and $T=10$, the admissible range of the damping $\alpha$, and
    Carr and Madan's choice of a quarter of its upper end.
\end{enumerate}
\end{exercise}

\begin{exercise}[Digitals and the skew]\label{ex:sv-digitals}
\medium A cash-or-nothing call pays $\1_{\{S_T>K\}}$; a gap call with trigger $K_1$ and strike $K_2$ pays
$(S_T-K_2)\1_{\{S_T>K_1\}}$.
\begin{enumerate}[label=(\alph*)]
  \item Price the gap call through $P_1$ and $P_2$ of \noteref{prop:sv-fourier} at the strike $K_1$, show that it is a
    call struck at $K_1$ plus $K_1-K_2$ cash-or-nothing calls, and evaluate it for $K_1=100$, $K_2=95$, $\rho=-0.7$.
  \item Compute the cash-or-nothing call at $K=100$ for $\rho=-0.7$ and $\rho=0$, and compare it with the
    Black--Scholes digital evaluated at the at-the-money implied volatility.
  \item Explain the difference with \noteref{prop:sk-binary-smile}, and check it numerically with a finite
    difference of implied volatilities.
\end{enumerate}
\end{exercise}

\begin{exercise}[Fixing Euler's scheme]\label{ex:sv-fixes}
\medium This is a computational exercise. Implement the four fixes of the notes (absorption, reflection,
partial and full truncation) with the log-Euler step for $S$.
\begin{enumerate}[label=(\alph*)]
  \item Show that, with the log-Euler step, every fix reproduces the forward exactly, and that the conditional
    estimator of \noteref{prop:sv-full-truncation}(ii) applies to every fix.
  \item Estimate the bias of the at-the-money put for $h=\frac14$ and $h=\frac1{16}$, in the test case with
    $\rho=-0.7$ and for $V_0=\theta=0.04$, $\kappa=0.5$, $\xi=1$, $\rho=-0.9$, against the Fourier prices
    $7.948022$ and $4.576954$. Report standard errors and seeds.
  \item Rank the fixes. Can you explain the ranking by comparing the drift that each fix applies when
    $\tilde V<0$?
\end{enumerate}
\end{exercise}

\begin{exercise}[Calibrating to a synthetic smile]\label{ex:sv-calibration}
\hard This is a computational exercise. You may use \texttt{scipy.optimize.least\_squares} and the pricers of the
chapter.
\begin{enumerate}[label=(\alph*)]
  \item Generate implied volatilities of out-of-the-money options from the test case with $\rho=-0.7$ for
    $K\in\{80,90,95,100,105,110,120\}$ and $T\in\{0.25,0.5,1,2\}$, and add independent normal noise with standard
    deviation $0.1\%$ (in volatility points), with a fixed seed.
  \item Calibrate $(V_0,\kappa,\theta,\xi,\rho)$ by least squares on implied volatilities from a start far from
    the truth. Report the estimates, the root-mean-square error, and approximate standard errors from the
    Jacobian. Which parameter is least well determined?
  \item Now use only the slice $T=1$, fix $\kappa\in\{0.5,2,8\}$ in turn and calibrate the other four. Compare the
    fits and the parameters, and the Feller ratio of each fit. What do you conclude about identifying the model
    from one maturity?
\end{enumerate}
\end{exercise}

\begin{exercise}[A formula without damping]\label{ex:sv-lewis}
\hard Let $k=\ln K$ and $\varphi$ as in the notes.
\begin{enumerate}[label=(\alph*)]
  \item Show that $\int_\R e^{ivz}(v^2+\frac14)^{-1}\,\dd v=2\pi e^{-|z|/2}$ for real $z$, and deduce that
    $\min(S_T,K)=\frac{\sqrt{S_TK}}{2\pi}\int_\R e^{iv(X_T-k)}(v^2+\frac14)^{-1}\,\dd v$.
  \item Deduce Lewis's (2001) formula
    \[
      C=Se^{-yT}-\frac{\sqrt K\,e^{-rT}}{\pi}\int_0^\infty\frac{\operatorname{Re}\big(e^{-ivk}\varphi(v-\frac i2)\big)}{v^2+\frac14}\,\dd v,
    \]
    and explain why, unlike Carr and Madan's, it needs no condition on the moments of $S_T$ beyond those that
    always hold.
  \item Obtain the same formula by moving the line of integration of Carr and Madan's \noteref{eq:sv-cm} from
    $\operatorname{Im}=-(1+\alpha)$ to $\operatorname{Im}=-\frac12$, and identify the residue you cross.
  \item Implement it with the midpoint rule on $[0,125]$ and check the put $7.948021924422$ of the notes
    ($\rho=-0.7$). How small must the step be, and why is it smaller than for Gil-Pelaez's formula?
\end{enumerate}
\end{exercise}
